\chapter{Bytes, encodings, and protocol layers}
\chapterintro{Networks move finite sequences of bits. Every line of text, JSON value, header field, and database result must eventually cross a boundary as bytes.}

\begin{objectives}
\begin{itemize}
  \item Distinguish bits, bytes, text, characters, and encodings.
  \item Explain framing and why a receiver needs an unambiguous boundary rule.
  \item Use layering and encapsulation without confusing the models with physical reality.
  \item Distinguish packets, TCP segments, byte streams, and application messages.
\end{itemize}
\end{objectives}

\section{A character is not a byte}

A \term{bit} has two possible values. Eight bits form a \term{byte}, an integer
from 0 through 255. A \term{character} is an abstract element of text. An
\term{encoding} maps characters to bytes and bytes back to characters.

ASCII assigns one byte to a small set of characters. UTF-8 represents every
Unicode code point using one to four bytes and deliberately matches ASCII for
the first 128 values. The word ``café'' has four characters but five UTF-8
bytes because \code{é} uses two. Lengths in an HTTP message describe bytes, not
Python characters.

\begin{mentalmodel}
Text does not travel through a socket. Bytes travel. Encoding is the contract
used before sending; decoding is the matching operation after receiving. A
protocol must state or imply that contract wherever text has meaning.
\end{mentalmodel}

In Python, \code{str} holds Unicode text and \code{bytes} holds byte values.
Combining them implicitly would hide an encoding decision, so Python requires an
explicit conversion. This apparent inconvenience prevents a large class of
wire-format mistakes.

\section{Framing creates messages}

Suppose a sender writes the bytes for \code{HELLO} followed immediately by the
bytes for \code{WORLD}. The receiver sees \code{HELLOWORLD}. Nothing inherent
in that stream identifies two words. A protocol needs a \term{framing} rule.

Common framing strategies include:

\begin{itemize}
  \item a delimiter, such as a newline ending a command;
  \item a length prefix, such as a four-byte integer followed by that many bytes;
  \item a fixed size, when every record occupies the same number of bytes;
  \item connection close, where end-of-stream terminates the message;
  \item self-describing chunks, each carrying its own length and a final marker.
\end{itemize}

HTTP/1.1 combines strategies. The header section uses a blank line delimiter.
The body length can come from method/status semantics, \code{Content-Length},
transfer coding, or closing a connection in limited circumstances. Correct
framing is security-critical: if two recipients choose different boundaries,
an attacker may make one request look like two.

\begin{wireview}
The notation below names bytes, not the four printable letters ``CRLF.''
\begin{wiremessage}
GET /health HTTP/1.1 CRLF
Host: localhost CRLF
CRLF
\end{wiremessage}
On the wire, each \code{CRLF} is the pair \code{0D 0A} in hexadecimal. The
second consecutive pair ends the header section.
\end{wireview}

\section{Layers manage complexity}

A layered model gives each protocol a focused job and a stable interface to the
layer above it. The practical Internet stack is often described with four or
five layers:

\begin{center}
\begin{tabularx}{\textwidth}{@{}>{\bfseries}p{30mm} p{32mm} X@{}}
\toprule
Layer & Examples & Responsibility \\
\midrule
Application & HTTP, DNS, TLS & Meaning and rules used by applications \\
Transport & TCP, UDP, QUIC & Communication between processes; reliability or datagrams \\
Internet & IPv4, IPv6 & Addressing and forwarding across networks \\
Link & Ethernet, Wi-Fi & Delivery on one local link \\
Physical & Copper, fiber, radio & Signals that represent bits \\
\bottomrule
\end{tabularx}
\end{center}

The seven-layer OSI model remains useful vocabulary, but Internet implementations
do not always fit its divisions exactly. TLS, for example, sits between an
application protocol and a transport in common usage. Models help reasoning;
they are not laws of nature.

\section{Encapsulation and changing envelopes}

When an application sends an HTTP request, each lower layer adds control
information. TCP groups some stream bytes into segments with sequence-related
metadata. IP carries a transport segment in a packet with source and destination
addresses. A link carries the packet in a frame suitable for the next hop.

\begin{figure}[htbp]
\centering
\begin{tikzpicture}[
  layer/.style={draw=Rule,rounded corners=1mm,minimum width=115mm,
                minimum height=10mm,align=center,font=\small\sffamily}
]
\node[layer,fill=PaperTint] (http) {HTTP message bytes};
\node[layer,fill=PaperTint,below=3mm of http,minimum width=130mm] (tcp)
  {TCP header \quad | \quad stream bytes};
\node[layer,fill=PaperTint,below=3mm of tcp,minimum width=145mm] (ip)
  {IP header \quad | \quad TCP segment};
\node[layer,fill=PaperTint,below=3mm of ip,minimum width=158mm] (link)
  {link header \quad | \quad IP packet \quad | \quad trailer};
\draw[-{Latex},Accent] (http) -- (tcp);
\draw[-{Latex},Accent] (tcp) -- (ip);
\draw[-{Latex},Accent] (ip) -- (link);
\end{tikzpicture}
\caption{Conceptual encapsulation. Sizes are illustrative, not proportional.}
\label{fig:encapsulation}
\end{figure}

At a router, the link envelope normally changes because the next link has new
local addressing. The IP packet continues toward its destination, though fields
such as the hop limit change. At the destination, layers remove and interpret
their own metadata before delivering stream bytes to the application.

\section{Packet boundaries are not message boundaries}

The word ``packet'' is often used loosely. Be precise when precision matters:
a link transmits frames, IP forwards packets, TCP transmits segments, and HTTP
defines messages. These boundaries need not align. One HTTP header can span
several TCP segments. Several short HTTP writes can be combined. One call to
receive bytes can return half a line, exactly one request, or a request plus the
beginning of the next request.

TCP preserves order and detects missing data; it does not preserve the sender's
write-call boundaries. Therefore a correct parser accumulates bytes and applies
HTTP's framing rules. It never assumes one receive operation equals one message.

\begin{pitfall}
Testing only on localhost often hides fragmentation bugs. Local transfers are
fast and may make a small request appear in one read. That outcome is permitted,
not guaranteed. Later tests deliberately feed one byte at a time.
\end{pitfall}

\section{Numbers and byte order}

Binary protocols must also agree how multi-byte numbers are ordered. Network
protocols traditionally use big-endian, often called network byte order: the
most significant byte comes first. HTTP/1.1's visible integers are decimal text,
so \code{Content-Length: 12} is the ASCII sequence for the characters 1 and 2,
not a binary integer. HTTP/2 and HTTP/3 use binary framing where byte order again
becomes visible to implementers.

\summarytitle

Bytes are the common currency of network interfaces. Encodings give bytes text
meaning; framing gives a stream message boundaries; layering gives complicated
systems manageable interfaces. TCP is an ordered stream, not a message queue,
so HTTP parsing must be incremental and boundary-aware.

\reviewtitle
\begin{enumerate}
  \item Why can four characters require more than four bytes?
  \item Name three framing strategies and one trade-off for each.
  \item Why may a link-layer frame change at every router?
  \item What outcomes are legal for a receive call after one HTTP request was sent?
\end{enumerate}

\begin{exercise}[Design a frame]
Design a tiny protocol for sending a sequence of UTF-8 notes. Choose a framing
strategy. Explain how the receiver detects an incomplete note, an empty note,
and the end of the sequence. Then identify one malformed input it must reject.
\end{exercise}

\begin{exercise}[Count the bytes]
For the texts \code{hello}, \code{café}, and a word in a language you use,
record character count and UTF-8 byte count. Explain which count belongs in a
length field that frames network bytes.
\end{exercise}
